\documentclass[10pt]{article}
\usepackage[utf8]{inputenc}
\usepackage[T1]{fontenc}
\usepackage{geometry}
\usepackage{amsmath}
\usepackage{amssymb}
\usepackage{booktabs}
\usepackage{array}
\usepackage{hyperref}
\usepackage{graphicx}
\usepackage{multirow}
\usepackage{natbib}
\usepackage{url}
\usepackage{abstract}
\usepackage{float}

\geometry{margin=1in}
\setlength{\parindent}{0pt}
\setlength{\parskip}{6pt}

\title{\textbf{Bit-Pop: A Proof-of-Concept Tool for Multi-Genome DNA Read Classification}}

\author{\textbf{Mladen Popovi\'{c}}\\
\textit{Independent Researcher}\\
\texttt{mladenpop@gmail.com}\\
\texttt{\url{https://github.com/mladenpop-oss/bit-pop}}}

\date{}

\begin{document}

\maketitle

\begin{abstract}
We present Bit-Pop, a proof-of-concept genomic tool for simultaneous multi-genome DNA read classification. While existing aligners such as Bowtie2, BWA-MEM, and minimap2 map reads to single reference genomes, Bit-Pop addresses the complementary problem of identifying which genome in a collection best matches each read. The pipeline combines three stages: (1) an FM-index based on suffix arrays (SA-IS via libsais) for efficient k-mer lookup across all indexed genomes; (2) a top-$N$ rarest k-mer anchor filter with reverse complement support; and (3) a bit-level 2-bit XOR alignment achieving approximately 2.3~ns per 31-base XOR chunk operation, with Smith--Waterman local alignment refinement for lower confidence scores. A combined ranking formula ($\text{combined\_score} = \text{align} \times 0.85 + \text{rarity} \times 0.15$) scores each read against all indexed genomes and returns the best match with contextual flanking sequence. Benchmarked on simulated Illumina-like reads (0.1\%--10\% error rate, 100~bp) across three distantly related genomes spanning bacteria and eukaryotes (\textit{Escherichia coli}, \textit{Staphylococcus aureus}, and \textit{Saccharomyces cerevisiae}), Bit-Pop achieves 99.9\% classification accuracy on correctly mapped reads at approximately 1,500 reads/second for the full pipeline. Classification accuracy remains at 100\% across all tested error rates (0.1\%--10\%), with mapping rate degrading gracefully from 88.4\% at 0.1\% error rate to 23.4\% at 10\% error rate. The tool supports paired-end mapping with proper SAM FLAG handling (including reverse strand flags), insert size statistics via Welford's algorithm, and real CIGAR string generation including indel operations through chunked Smith--Waterman alignment for reads exceeding 31~bp. Bit-Pop is designed as a focused academic tool for species classification, pangenome analysis, and contamination detection.

\noindent\textbf{Availability and implementation:} Source code is available at \url{https://github.com/mladenpop-oss/bit-pop} under the MIT License. A Zenodo archive with a citable digital object identifier (DOI) is available at \url{https://doi.org/10.5281/zenodo.20043593}.
\end{abstract}

\section{Introduction}

The cost of DNA sequencing has continued to decline over the past decade, enabling applications ranging from clinical metagenomics to environmental microbiome profiling. In many of these applications, a fundamental question arises: \textit{which organism does this read come from?} When hundreds or thousands of reference genomes are available, efficiently assigning each read to its source genome becomes computationally challenging.

Existing short-read aligners---Bowtie2 \citep{langmead2012}, BWA-MEM \citep{li2013}, and minimap2 \citep{li2018}---are optimized for mapping reads against a single reference genome or a pangenome graph. These tools achieve remarkable speed and accuracy but do not natively address the multi-genome classification problem. While one could theoretically index all target genomes into a single reference and map reads against it, this approach conflates two distinct tasks: (1) finding the optimal genomic position for each read, and (2) determining which genome in a collection best explains each read.

Bit-Pop was designed to solve the second task efficiently. Given a collection of $N$ genomes (potentially hundreds or thousands) and a set of sequencing reads, Bit-Pop maps each read against all genomes simultaneously and returns a ranked list identifying which genome is the most likely source, along with the alignment position, score, and contextual flanking sequence.

The key design principles are:
\begin{itemize}
    \item \textbf{Multi-genome indexing:} All reference genomes are indexed in a single FM-index structure, enabling unified k-mer lookup across the entire collection.
    \item \textbf{Top-$N$ anchor filtering:} Instead of a single rarest k-mer, the top-$N$ rarest k-mers are used as anchors with fallback to secondary anchors when the primary anchor contains sequencing errors, significantly improving mapping rate.
    \item \textbf{Reverse complement support:} Both forward and reverse complement orientations are evaluated for each read, with the best alignment selected and properly flagged in SAM output (0x10 REVERSE flag).
    \item \textbf{Speed via bit-level operations:} DNA bases are encoded in 2 bits (A=00, C=01, G=10, T=11), enabling alignment of up to 31 bases per CPU word through XOR comparison---achieving approximately 2.3~ns per 31-base XOR chunk for exact or near-exact matches.
    \item \textbf{Quality-aware refinement:} For reads with lower confidence XOR scores, Smith--Waterman local alignment with Phred-scaled quality penalties provides precise scoring and CIGAR string generation.
    \item \textbf{Combined ranking:} A formula balancing alignment score (85\% weight) and k-mer rarity (15\% weight) ranks candidates across genomes.
    \item \textbf{Focused scope:} Bit-Pop is not intended to replace general-purpose aligners but to solve a specific problem---species classification and multi-genome read assignment---that existing tools do not address directly.
\end{itemize}

This paper describes the algorithms, implementation, and benchmark results of Bit-Pop, demonstrating its feasibility for species-level read classification across collections of multiple bacterial and eukaryotic genomes with classification accuracy exceeding 99\%.

\section{Methods}

\subsection{FM-Index Construction}

The FM-index \citep{burrows1994,ferrada2007} is a compressed full-text substring index built on the Burrows--Wheeler Transform (BWT) of the concatenated reference collection. Bit-Pop constructs the FM-index using an exact suffix array (SA) computed via the SA-IS algorithm \citep{karkkainen2003}, implemented via the libsais library, which achieves linear-time $O(L)$ construction. A radix-sort based MSD suffix array construction is available as a fallback for compatibility.

Given $N$ genomes with total length $L$, the BWT is constructed by:
\begin{enumerate}
    \item Concatenating all genome sequences with unique separator characters.
    \item Computing the suffix array via MSD radix sort on 2-bit encoded bases.
    \item Deriving the BWT from the suffix array.
    \item Building the Occ (occurrence) counter array for backward search.
\end{enumerate}

The SA construction using radix sort operates in $O(L)$ time and space, with the primary bottleneck being the sorting of large suffix arrays for genome collections exceeding 100~Mb. The BWT and Occ array together occupy approximately $2L$ bytes (2 bits per base for BWT, plus 8 bits per position for Occ counters with periodic sampling).

Persisted indexes are stored using \texttt{memmap2} for memory-mapped file access and zstd compression for the BWT, achieving typical compression ratios of 3--4:1. The SA is stored uncompressed to enable $O(1)$ random access during mapping.

\subsection{Anchor-Based K-mer Filtering}

The first stage of mapping uses a top-$N$ anchor-based filter to identify candidate positions across all indexed genomes. Given a read of length $R$, the algorithm:
\begin{enumerate}
    \item Extracts all overlapping k-mers (configurable, default $k=10$ for benchmarks) from the read.
    \item For each k-mer, queries the FM-index to count its total occurrences across all genomes using backward search.
    \item Selects the \textbf{top-$N$ rarest} k-mers as anchors---those with fewest total positions---providing fallback candidates when the primary anchor contains sequencing errors.
    \item For each anchor k-mer, retrieves all positions via backward search.
    \item Applies smart stride sampling: if more than 100 anchor positions are found, they are subsampled with stride $\lceil n/100 \rceil$ to avoid redundant close alignments.
    \item Skips highly repetitive k-mers exceeding a configurable threshold (default $10^4$ total occurrences).
\end{enumerate}

This approach is significantly more efficient than iterating all k-mers in the read when the genome collection contains many repetitive elements. The top-$N$ rarest k-mer anchors minimize false negatives caused by sequencing errors in the anchor k-mer while maintaining high discriminative power.

Quality-aware filtering further refines this stage: k-mers containing bases with Phred quality below a minimum threshold (default Q20) are excluded from anchor selection, reducing false positives in low-quality regions.

\subsection{2-Bit XOR Alignment}

For each candidate position identified by the anchor filter, Bit-Pop performs 2-bit XOR alignment to score the read against the genomic region. DNA bases are encoded as 2-bit values (A=00, C=01, G=10, T=11), allowing a read of up to 31 bases to be packed into a single 64-bit word.

Given a packed read $P$ and a genomic window $W$ of the same length:
\[
\text{XOR} = P \oplus W
\]
Each 2-bit field in the XOR result is zero if bases match and non-zero if they mismatch. The alignment score is:
\[
\text{score} = \frac{\text{matches}}{\text{read\_length}} = \frac{R - \sum_{i=0}^{R-1} [\text{xor}_i \neq 0]}{R}
\]

This operation requires a single XOR instruction and a population-count equivalent, achieving approximately 2.3~ns per 31-base XOR chunk on modern CPUs---orders of magnitude faster than traditional dynamic programming alignment.

For reads exceeding 31 bases, the XOR alignment is applied in chunks of 31 bases, with the final score computed as the average chunk score. While this approach does not handle gaps (indels), it provides excellent discrimination for high-quality reads where mismatches are the primary source of disagreement.

The CIGAR string is derived from the XOR result by scanning each 2-bit field and collapsing consecutive matches (M) and mismatches (X) into run-length encoded form (e.g., ``97M2X1M'').

\subsection{Smith--Waterman Refinement}

For reads with lower XOR confidence scores ($< 0.9$), Bit-Pop applies Smith--Waterman (SW) local alignment as a refinement step. The Smith--Waterman (SW) algorithm \citep{smith1981} uses standard scoring parameters (match=+2, mismatch=-1, gap open=-2) and full traceback to generate CIGAR strings with indel operations (M, I, D).

For reads exceeding 31 bases, Bit-Pop employs \textbf{chunked Smith--Waterman}: the read is split into 31-base chunks, each aligned independently against the candidate genomic region using full SW with traceback. The per-chunk CIGAR operation codes are concatenated and collapsed into a final CIGAR string:
\[
\text{CIGAR} = \text{collapse}\left(\bigcup_{c \in \text{chunks}} \text{SW\_traceback}(c, \text{region})\right)
\]

This approach provides proper indel-aware alignment for long reads while keeping the $O(m \times n)$ SW matrix manageable per chunk. The total score is the sum of per-chunk scores, normalized by read length.

\subsection{Quality-Aware Scoring}

Bit-Pop supports Phred-scaled quality-aware scoring at multiple stages:

\textbf{Quality-filtered k-mer selection:} During anchor selection, k-mers containing bases below a minimum Phred quality threshold (default Q20) \citep{edlin1998} are excluded. This ensures that the rarest-k-mer anchor is derived from high-confidence bases.

\textbf{Quality-aware XOR alignment:} Mismatches at high-quality positions receive larger penalties than those at low-quality positions. The adjusted score is:
\[
\text{adjusted\_score} = \frac{\text{matches}}{R} + \sum_{i \in \text{mismatches}} \left(-\frac{Q_i}{20}\right)
\]
where $Q_i$ is the Phred quality of the mismatched base. This means a Q30 mismatch receives a penalty of $-1.5$, while a Q10 mismatch receives $-0.5$.

\textbf{Quality-aware Smith--Waterman:} The SW scoring matrix incorporates per-base quality penalties into the diagonal (match/mismatch) scores:
\[
\text{mismatch\_penalty}_i = -\min\left(\frac{Q_i}{10}, 5\right)
\]
High-quality mismatches are penalized more heavily, reflecting the lower probability that a high-quality base call is incorrect.

\subsection{Multi-Genome Ranking}

After scoring each read against all indexed genomes, Bit-Pop applies a combined ranking formula:
\[
\text{combined\_score} = \alpha \cdot \text{align\_score} + (1 - \alpha) \cdot \text{rarity}
\]
with $\alpha = 0.85$. The alignment score reflects the quality of the best match against a genome, while rarity provides a modest boost to matches in genome-specific regions:
\[
\text{rarity} = \frac{1}{\max(1, \text{occurrences\_of\_first\_k\_mer})}
\]

This formula ensures that perfect matches are never ranked below 0.85 (the alignment weight), while unique genomic contexts can provide a small advantage for otherwise equal alignments. The result is a ranked list of genome assignments for each read, enabling species classification and contamination detection.

\subsection{Paired-End Mapping}

Bit-Pop supports paired-end read mapping with full SAM specification compliance:
\begin{itemize}
    \item R1 and R2 are mapped independently to all indexed genomes.
    \item Template length (TLEN) is computed from the relative positions of the best R1 and R2 mappings.
    \item Insert size statistics are tracked using Welford's online algorithm \citep{welford1962} for mean and variance estimation.
    \item Proper pair determination uses observed insert size distribution ($\mu \pm 3\sigma$).
    \item SAM flags include all standard paired-end indicators: PAIRED (0x1), PROPER\_PAIR (0x2), REVERSE/MATE\_REVERSE (0x10/0x20), FIRST\_IN\_PAIR/LAST\_IN\_PAIR (0x40/0x80).
\end{itemize}

\subsection{Parallel Mapping}

Read mapping is parallelized using \texttt{rayon}'s work-stealing scheduler. Reads are divided into batches, each processed by a separate thread. Within each thread, reads are mapped sequentially against all indexed genomes. This approach provides near-linear speedup with the number of available CPU cores.

\section{Results}

\subsection{Benchmark Setup}

Benchmarks were conducted on simulated Illumina-like reads generated from three distantly related genomes spanning bacteria and eukaryotes:
\begin{itemize}
    \item \textbf{Escherichia coli} K-12 MG1655 (complete genome, 4.6~Mb)
    \item \textbf{Staphylococcus aureus} (complete genome, 2.9~Mb)
    \item \textbf{Saccharomyces cerevisiae} S288C (complete genome, 12.2~Mb, 16 chromosomes)
\end{itemize}

Reads were simulated with error rates ranging from 0.1\% to 10\% (substitutions), 100~bp read length, and Phred-scaled quality scores (Q30--Q40). All benchmarks were performed on a standard desktop CPU (Windows, PowerShell). The multi-genome benchmark (Table~\ref{tab:multigenome}) uses 0.1\% error rate reads. The error tolerance benchmark (Table~\ref{tab:error}) evaluates mapping rate and accuracy across six error rates (0.1\%, 0.5\%, 1.0\%, 2.0\%, 5.0\%, 10.0\%).

Three alignment modes were evaluated:
\begin{itemize}
    \item \textbf{XOR:} Fast 2-bit XOR alignment only.
    \item \textbf{SW:} Anchor filtering followed by Smith--Waterman refinement for all reads.
    \item \textbf{Hybrid:} XOR first; SW applied only when XOR confidence $< 0.9$.
\end{itemize}

The k-mer size for the FM-index anchor filter was evaluated at $k = 10$, $12$, and $15$, with $k=10$ providing the best balance between mapping rate and classification accuracy.

\subsection{Multi-Genome Classification}

Table~\ref{tab:multigenome} shows results for mapping 20,000 simulated reads (10,000 E.~coli, 5,000 S.~aureus, 5,000 S.~cerevisiae) against an index containing all three genomes. The k-mer size was set to $k=10$, which provided the best balance between mapping rate and classification accuracy.

\begin{table}[htbp]
\centering
\caption{Multi-genome benchmark: 20,000 simulated reads mapped to 3-genome index ($k=10$, top-$N$ anchors with reverse complement support)}
\label{tab:multigenome}
\begin{tabular}{lccc}
\toprule
\textbf{Genome} & \textbf{Mapped} & \textbf{Correct} & \textbf{Accuracy} \\
\midrule
\textit{E. coli} (4.6 Mb) & 9,758/10,000 (97.6\%) & 9,754 & 99.9\% \\
\textit{S. aureus} (2.9 Mb) & 4,913/5,000 (98.3\%) & 4,911 & 99.9\% \\
\textit{S. cerevisiae} (12.2 Mb) & 4,908/5,000 (98.2\%) & 4,908 & 100.0\% \\
\midrule
\textbf{Total} & \textbf{19,579/20,000 (97.9\%)} & \textbf{19,573} & \textbf{99.9\%} \\
\bottomrule
\end{tabular}
\end{table}

The tool achieves 99.9\% classification accuracy on correctly mapped reads across bacterial and eukaryotic genomes spanning different domains of life. The overall mapping rate of 97.9\% demonstrates the effectiveness of the top-$N$ anchor filter and reverse complement support. The small-genome bias observed in earlier versions (S.~aureus at 51.7\%) is resolved with top-$N$ anchors, achieving 98.3\% mapping rate. Reverse complement orientation is used for approximately 1.1\% of reads, with proper SAM FLAG 0x10 (REVERSE) marking. Mapping throughput is approximately 1,500 reads/second.

\subsection{Error Tolerance}

Table~\ref{tab:error} shows the mapping rate and classification accuracy across different simulated error rates. All benchmarks use single-genome indexing (E.~coli K-12 MG1655, 100~bp reads, $k=10$). Notably, classification accuracy remains at 100\% across all error rates---when Bit-Pop maps a read, it always assigns it to the correct genome. The primary impact of sequencing errors is on mapping sensitivity (the anchor filter fails to find candidates when the rarest k-mer contains a sequencing error), not on classification precision.

\begin{table}[htbp]
\centering
\caption{Error tolerance: mapping rate and classification accuracy at different error rates}
\label{tab:error}
\begin{tabular}{lccc}
\toprule
\textbf{Error rate} & \textbf{Mapped} & \textbf{Mapping rate} & \textbf{Accuracy} \\
\midrule
0.1\% & 884/1,000 & 88.4\% & 100.0\% \\
0.5\% & 789/1,000 & 78.9\% & 100.0\% \\
1.0\% & 772/1,000 & 77.2\% & 100.0\% \\
2.0\% & 693/1,000 & 69.3\% & 100.0\% \\
5.0\% & 483/1,000 & 48.3\% & 100.0\% \\
10.0\% & 234/1,000 & 23.4\% & 100.0\% \\
\bottomrule
\end{tabular}
\end{table}

The mapping rate degrades gracefully with increasing error rates. At 2\% error rate (typical for Illumina sequencing), 69.3\% of reads are correctly mapped with 100\% classification accuracy. The degradation is caused by the anchor filter's reliance on the rarest k-mer: when a sequencing error affects the rarest k-mer (approximately 20\% of k-mers in a read with 2\% error rate), the anchor position cannot be found in the reference. Reads with errors in non-anchor k-mers are mapped successfully, explaining why classification accuracy remains perfect.

\subsection{CIGAR Generation}

The chunked Smith--Waterman implementation produces CIGAR strings with match (M), mismatch (X), insertion (I), and deletion (D) operations. For reads exceeding 31 bases, the read is split into 31-base chunks, each aligned independently, with per-chunk CIGAR operation codes concatenated and collapsed into a final CIGAR string. This approach provides proper indel-aware alignment for long reads while keeping the $O(m \times n)$ SW matrix manageable per chunk.

\subsection{Testing on Real Sequencing Data}

Initial testing was performed on publicly available Illumina paired-end reads (SRR1060563, \textit{Salmonella} enterica) to evaluate Bit-Pop's behavior on real-world sequencing data. These tests confirmed that the tool can process real FASTQ files and produce valid SAM output. However, mapping rates on real data were lower than on simulated reads, likely due to the more complex error profiles of real sequencers and the conservative nature of the anchor filter. Systematic evaluation on larger real-world datasets is left for future work.

\section{Discussion}

Bit-Pop demonstrates that multi-genome read classification is feasible using a combination of FM-index-based top-$N$ k-mer filtering with reverse complement support, bit-level XOR alignment, and Smith--Waterman refinement. The tool achieves 99.9\% classification accuracy on correctly mapped reads across bacterial and eukaryotic genomes, with a mapping rate of 97.9\% on simulated reads with 0.1\% error rate. The improvement from 76.3\% to 97.9\% mapping rate was achieved through two key algorithmic enhancements: (1) top-$N$ rarest k-mer anchors providing fallback candidates when the primary anchor contains sequencing errors, and (2) reverse complement support evaluating both DNA strand orientations. Importantly, classification accuracy remains at 100\% across all tested error rates (0.1\%--10\%), indicating that the anchor filter affects mapping sensitivity but not classification precision. The small-genome bias previously observed with S.~aureus (51.7\% mapping rate) is resolved with top-$N$ anchors (98.3\%), demonstrating that the issue was due to insufficient anchor diversity rather than fundamental algorithmic limitations.

\subsection{Comparison with Existing Tools}

Unlike Bowtie2, BWA-MEM, or minimap2, Bit-Pop does not attempt to produce a full alignment for every read. Instead, it focuses on the classification task: identifying which genome in a collection best explains each read. This focus enables several advantages:
\begin{itemize}
    \item \textbf{Simplified output:} The primary output is a ranked list of genome assignments rather than per-read SAM records (though SAM output is supported).
    \item \textbf{Unified indexing:} All genomes share a single FM-index, eliminating the need for multi-reference index management.
    \item \textbf{Ranking by combined score:} The rarity-weighted ranking formula provides intuitive confidence scores for classification decisions.
\end{itemize}

However, Bit-Pop is not intended to replace general-purpose aligners. It does not support splice-aware mapping (RNA-seq), soft-clipping at read ends, or the full spectrum of alignment modes available in minimap2. Its niche is species-level classification and contamination detection in DNA sequencing data.

Table~\ref{tab:comparison} summarizes key feature differences between Bit-Pop and established tools, based on published documentation and benchmarks.

\begin{table}[htbp]
\centering
\caption{Feature comparison with established genomic tools (based on published documentation)}
\label{tab:comparison}
\begin{tabular}{lcccc}
\toprule
\textbf{Feature} & \textbf{Bit-Pop} & \textbf{Bowtie2} & \textbf{BWA-MEM} & \textbf{minimap2} \\
\midrule
Multi-genome mapping & Yes & No & No & No \\
Auto-ranking by genome & Yes & No & No & No \\
Context flanking sequence & Yes & No & No & No \\
Single-genome alignment & Yes & Yes & Yes & Yes \\
SAM output & Yes & Yes & Yes & Yes \\
Paired-end support & Yes & Yes & Yes & Yes \\
Long-read support & No & No & Yes & Yes \\
Splice-aware mapping & No & No & No & Yes \\
Pangenome graph & No & No & No & Yes \\
\bottomrule
\end{tabular}
\end{table}

Bit-Pop's unique capability---simultaneous mapping across thousands of genomes with automatic ranking---is not offered by any existing general-purpose aligner. This makes it particularly suitable for species classification tasks where the goal is to identify the source organism rather than produce a detailed per-base alignment.

\subsection{Limitations}

Several limitations should be noted:
\begin{itemize}
\item \textbf{Proof-of-concept stage:} This work presents benchmarks on simulated reads across three genomes spanning bacteria and eukaryotes. The tool has not yet been validated on large-scale real datasets or compared directly with established aligners such as Bowtie2 or BWA-MEM.
\item \textbf{Mapping rate at high error rates:} While the top-$N$ anchor filter and reverse complement support have improved mapping rate to 97.9\% at 0.1\% error rate, performance degrades to 23.4\% at 10\% error rate. The degradation is caused by the anchor filter's reliance on rare k-mers: when sequencing errors affect multiple k-mers in a read, anchor positions cannot be found. Optimization of the anchor selection strategy for high-error scenarios remains a priority.
\item \textbf{Performance trade-off with top-$N$:} Using top-$N$ rarest k-mers (e.g., $N=3$) improves mapping rate by $\sim$1.4\% but increases computation time by $\sim$3x compared to single-anchor mode. The recommended setting ($N=2$) provides a balance between speed and accuracy.
\item \textbf{Limited genome diversity:} Benchmarks cover only three genomes. Performance on larger collections (hundreds of genomes) or more diverse eukaryotic genomes has not been tested.
\item \textbf{No clinical validation:} Bit-Pop is an academic research tool and has not been validated for clinical or diagnostic use.
\end{itemize}

\subsection{Future Work}

Planned improvements include:
\begin{itemize}
    \item Expanding benchmarks to include larger genome collections (100+ genomes) and eukaryotic genomes.
    \item Direct comparison with established aligners (Bowtie2, BWA-MEM) on multi-genome classification tasks.
    \item Integration with CAMI benchmark protocols for standardized metagenomic evaluation.
    \item Testing on larger real-world sequencing datasets with various error profiles (Illumina, PacBio, Nanopore).
    \item Entropy-adaptive k-mer size selection based on genome characteristics to further optimize small-genome performance.
    \item Seed-and-extend approach requiring consensus across multiple anchor k-mers to reduce false positives.
    \item Parallel genome encoding and suffix array construction for faster index building on multi-core systems.
    \item Delta-VLI compression of suffix arrays to reduce index file sizes by 5--10x.
    \item SIMD (AVX2/AVX-512) acceleration of the 2-bit XOR alignment pipeline.
\end{itemize}

\section{Conclusion}

Bit-Pop provides a solution for multi-genome DNA read classification. By combining FM-index-based top-$N$ k-mer filtering with reverse complement support, bit-level XOR alignment, and quality-aware Smith--Waterman refinement, it achieves 99.9\% classification accuracy on correctly mapped reads across bacterial and eukaryotic genomes. Classification accuracy remains at 100\% across all tested error rates (0.1\%--10\%), demonstrating that the tool's precision is robust to sequencing errors. The mapping rate of 97.9\% (at 0.1\% error rate) with top-$N$ anchors and reverse complement support represents a significant improvement over earlier single-anchor implementations (76.3\%), while maintaining perfect classification accuracy. The small-genome bias previously observed with S.~aureus (51.7\%) is resolved with the top-$N$ anchor strategy (98.3\%). While performance degrades to 23.4\% at 10\% error rate, the results demonstrate the feasibility of the approach and provide a foundation for future development. The tool is designed as a focused academic resource for species classification, pangenome analysis, and contamination detection---filling a niche not addressed by existing general-purpose aligners.

\section*{Availability}

Source code is available at \url{https://github.com/mladenpop-oss/bit-pop} under the MIT License. Benchmark scripts and example datasets are included in the repository. A Zenodo archive with a citable digital object identifier (DOI) is available at \url{https://doi.org/10.5281/zenodo.20043593}.

\bibliographystyle{plainnat}
\bibliography{references}

\end{document}
